\chapter{Conclusioni}
\label{cha:conclusions}

I risultati delle esecuzioni mostrano risultati chiari e confortanti, motori come QLever, QEndpoint e Virtuoso riescono ad avere
prestazioni elevate anche su questi carichi di lavoro, non è quindi preclusa la possibilità che Big Data di diverso tipo
vengano rappresentati attraverso RDF e poi analizzati su calcolatori non professionali, che siano quindi alla portata di tutti.
\par Va di certo fatto presente come questo non sia interalmente possibile né semplice al giorno d'oggi. La costruzione
del visualizzatore per il framework 3Dont ha richiesto molta ricerca in ambito di triple store, parametri di configurazione
e tecnologie di integrazione. Motori come QLever mettono a disposizione delle interfaccie native sotto forma di libreria,
ma nonostante questo non hanno ancora raggiunto livelli di usabilità ottimali, basti pensare che ad esempio anche nella sua 
forma integrata la libreria di QLever restituisce dalle query SPARQL soltanto dati serializzati, che vanno quindi deserializzati comunque. 
L'interfaccia HTTP di SPARQL rimane quindi molto pratica per la prototipazione ma limitante per il raggiungimento di
livelli di esperienza utente desiderabili, e richiede l'esecuzione del triple store in un processo dedicato. Processo che
in molti casi, soprattutto con soluzioni commerciali, è disponibile solo attraverso un container OCI, un ulteriore strato
di complessità e latenza che non rende di certo invitante l'integrazione.
\par Fanno ben sperare però i progressi e l'indirizzazione dello sviluppo di molti motori degli ultimi anni. I progetti più recenti
analizzati sono Oxigraph, QEndpoint e QLever, che insieme mostrano un interesse a raggiungere buone prestazioni in query di tipo
OLAP e con grandi scansioni dei dati, oltre che alla facilità di integrazione come librerie preconfezionate.


% TODO: gli strumenti analizzati non sono ancora utilizzabili in autonomia, servono
% integrazioni per l'accesso diretto ai modelli.
\section{Sviluppi Futuri}
\label{sec:future_dev}

Come premesso, questa tesi non vuole essere nulla di più di una breve analisi qualitativa e quantitativa della possibilità
di usare RDF per l'immagazzinamento di Big Data, sarebbe molto utile procedere con questa idea individuando dei casi d'uso
complessi e identificando i pattern di complessità ricorrenti oppure marginali ma comunque impattanti.
\par Va inoltre fatto notare che le tecnologie utilizzate come riferimento qui, RDF, Owl, SPARQL e l'ecosistema del semantic web,
sono solo uno delle piattaforme per lo sviluppo di software basato sulla conoscenza semantica e sui linked data. Esistono
diverse alternative come i Labeled Property Graphs (LPGs) che usano linguaggi come Cypher o Gremlin per l'interrogazione,
oppure il Property Graph Exchange Format (PGXF) con il Graph Query Language (GQL), standardizzato a livello globale dalla ISO/IEC.


\subsection{Intelligenza artificiale agentica}
\label{sec:agentic_ai}

Di fronte ai grandi progressi dell'intelligenza artificiale agentica degli ultimi anni diventa interessante analizzare
le potenzialità che l'uso di RDF per l'immagazzinamento e l'analisi di Big Data ci prospetta in questo campo.
Attualmente infatti lo standard de-facto usato dai modelli IA agentici per interagire con componenti esterni è il
framework Model Context Protocol (MCP). Questo framework permette di definire un connettore tra un software e un modello
IA agentico, ogni software diverso richiede la scrittura di un proprio componente. Usando un solo componente per interfacciarsi
con un grafo RDF o con un endpoint SPARQL sarebbe quindi possibile andare ad interagire con un vasto numero di formati di dato
differenti. Un agente IA è perfettamente in grado di comprendere un'ontologia autonomamente e l'espressività e idiomacità
del framework RDF si sposa sorprendentemente bene con le capacità dei moderni LLM.
\par Per condurre una reale valutazione delle potenzialità di questo metodo sarebbe necessario istruire un agente in 
modo minimo sull'ontologia che verrà utilizzata per poi richiederne l'analisi e l'esecuzione di un dato compito, per poi
confrontare i risultati con uno stesso agente IA che però è andato ad usare un'interfaccia MCP diretta per i dati analizzati.
Se i risultati fossero soddisfacienti si aprirebbero le porte all'automatizzazione delle operazioni su un gran numero
di dati disparati. todo
% TODO (sviluppi futuri): molto rilevante con le recenti piattaforme e strumenti di AI
% agentica, permette un'ampia rappresentazione della conoscenza e un accesso semplice
% tramite interfaccia standard (es. mappatura di codebase).
% https://github.com/Egonex-AI/Understand-Anything
